\chapter{Fermion Doubling Problem in the ``naive'' action}

The problem with the naive discretization of the Dirac action, specified
by the operator $M^N$ in Eq.~7.7, is that in the continuum limit it gives
rise to $2^d = 16$ flavors rather than one. This doubling problem is readily
demonstrated by the inverse of the free field propagator (obtained by taking
the fourier transform of the action with all $U_\mu(x) = 1$).
\begin{equation}
S^{-1}(p) = m_q + \frac{i}{a}\sum_\mu \gamma_\mu \sin p_\mu a
\end{equation}
which has 16 zeros within the Brillouin cell in the limit $m_q \to 0$. Defining
the momentum range of the Brillouin cell to be $\{-\pi/2, 3\pi/2\}$, the zeros
lie at $p_\mu = 0$ and $\pi$. As discussed later, this proliferation holds under very
general conditions specified by the Nielsen-Ninomiya theorem [14], and
is intimately related to simultaneously preserving chiral symmetry. The
inclusion of the gauge fields does not solve the doubling problem.

Before discussing partial fixes, it is useful to investigate the properties
of these extra zero modes under chiral transformations. Let us define a
set of 16 4-vectors $\Pi^A = \{(0, 0, 0, 0), \allowbreak (\pi, 0, 0, 0), \allowbreak \ldots, \allowbreak (\pi, \pi, \pi, \pi)\}$ with
$A = \{1 \ldots 16\}$, and consider the expansion of the massless propagator about
these points. Then
\begin{equation}
\begin{aligned}
S^{-1}(p, m = 0) &= \frac{i}{a}\sum_\mu \gamma_\mu \sin p_\mu a \\
&= \frac{i}{a}\sum_\mu \gamma_\mu \sin(\Pi^A + k)_\mu a \\
&= \frac{i}{a}\sum_\mu \gamma_\mu S_\mu^A \sin k_\mu a \\
&\equiv \frac{i}{a}\sum_\mu \tilde{\gamma}_\mu \sin k_\mu a ,
\end{aligned}
\end{equation}
where $S_\mu^A = \{+1, -1\}$ depending on whether the $\mu$ component of $\Pi^A$ is
$\{0, \pi\}$. In the last expression a new representation of gamma matrices has
been defined by the similarity transformation $\tilde{\gamma}_\mu = Y_A \gamma_\mu Y_A^\dagger$ where
\begin{equation}
Y_A = \prod_\mu (\gamma_\mu \gamma_5)^{n_\mu^A}
\end{equation}
and $n_\mu^A$ are $\{0, 1\}$ depending on whether the momentum expansion in the
$\mu$ direction is about 0 or $\pi$. Now, $\tilde{\gamma}_5 = Y_A \gamma_5 Y_A^\dagger =
S_1^A S_2^A S_3^A S_4^A \gamma_5 \equiv X^A \gamma_5$, with the signs $X^A$ given in table 3.
Thus, the sixteen species break up into two sets of 8 with chiral charge
$\pm 1$, and render the theory anomaly-free. Consequently, this ``naive''
discretization is phenomenologically not acceptable. The presence of
doublers can, in fact, be traced to a larger flavor symmetry of the naive
fermion action as I now discuss.

\begin{table}[htbp]
\centering
\begin{tabular}{ccc}
\toprule
$A$ & Degeneracy & $X^A$ \\
\midrule
$(0, 0, 0, 0)$ & 1 & $+1$ \\
$(1, 0, 0, 0)$ & 4 & $-1$ \\
$(1, 1, 0, 0)$ & 6 & $+1$ \\
$(1, 1, 1, 0)$ & 4 & $-1$ \\
$(1, 1, 1, 1)$ & 1 & $+1$ \\
\bottomrule
\end{tabular}
\caption{The chiral sign and degeneracy factor for the doublers.}
\label{tab:3}
\end{table}

In the continuum the mass-less Dirac action is invariant under the
$U(1)_V \otimes U(1)_A$ flavor transformations
\begin{equation}
\begin{aligned}
\psi_L &\to V_L \psi_L , & \bar{\psi}_L &\to \bar{\psi}_L V_L^\dagger , \\
\psi_R &\to V_R \psi_R , & \bar{\psi}_R &\to \bar{\psi}_R V_R^\dagger .
\end{aligned}
\end{equation}
Since the left and right handed fields are only coupled through the mass
term, the vector symmetry $V_L = V_R = e^{i\theta}$ holds for all $m$, whereas the
axial symmetry $V_L = V_R^\dagger = e^{i\theta}$ (or equivalently $V_L = V_R = e^{i\theta\gamma_5}$) holds
only for $m = 0$. The naive fermion action $S_N$ for a single flavor has a much
larger symmetry group, $U(4)_V \otimes U(4)_A$, under
\begin{equation}
\begin{aligned}
\psi(x) &\to \Gamma_x \left[ \frac{1 - \epsilon(x)\gamma_5}{2} V_L + \frac{1 + \epsilon(x)\gamma_5}{2} V_R \right] \Gamma_x^\dagger \psi(x) \\
\bar{\psi}(x) &\to \bar{\psi}(x) \Gamma_x \left[ V_L^\dagger \frac{1 + \epsilon(x)\gamma_5}{2} + V_R^\dagger \frac{1 - \epsilon(x)\gamma_5}{2} \right] \Gamma_x^\dagger
\end{aligned}
\end{equation}
if one considers translations by $2a$. $\Gamma_x$ is given in Eq.~11.1, the phase
$\epsilon(x)$ in Eq.~11.4, and $\Gamma_x \gamma_5 \Gamma_x^\dagger = \epsilon(x)\gamma_5$.
Here $V_{L,R}$ are general $4 \times 4$ unitary matrices acting on the spinor
index of the fermions. Even though there is, in Eq.~9.5, a dependence on $x$
through $\epsilon(x)$, this is a global symmetry as specifying it at one point
fixes it at all points.

To see how doubling arises, consider the vector ($V_L = V_R$) discrete
subgroup with the thirty two elements $\{\pm Y_A\}$ defined in Eq.~9.3. Eq.~9.2
shows that the $Y_A$ shift the momenta by $\pi/a$ along the direction specified
by $A$,
\begin{equation}
Y_A S_F(k) Y_A^\dagger = S_F\!\left( k + \frac{\pi}{a} A \right) .
\end{equation}
Since this is a similarity transform, the sixteen regions of the Brillouin
zone related by this transormation are physically equivalent. This analysis,
which relied only on the spinor structure of the action, makes it easy to
see why the doublers exist even for the interacting theory. In Section 11
this symmetry will be exploited to construct staggered fermions. In this
formulation the sixteen doublers are reduced to four, and at the same time
a $U(1)_V \otimes U(1)_\epsilon$ symmetry (corresponding to when $V_{L,R}$ are just phases)
is retained. The $U(1)_\epsilon$ plays the role of a chiral symmetry at finite $a$ and
in the continuum limit becomes the $U(1)_A$.

\section{Generalities of Chiral Symmetry}

In the massless limit of the continuum theory, the flavor symmetry at the
classical level is $U_V(1) \times U_A(1) \times SU_V(n_f) \times SU_A(n_f)$, where $n_f$ is the
number of light fermions. The unbroken $U_V(1)$ gives rise to baryon number
conservation. The $U_A(1)$ is broken by instanton contributions and the flavor
singlet axial current is anomalous. The $SU_A(n_f)$ is spontaneously broken
and the quark condensates acquire a non-zero expectation value in the
QCD vacuum. Associated with this spontaneously broken symmetry are
$(n_f^2 - 1)$ goldstone bosons, i.e. the pions. Chiral symmetry has played a
very useful role in the continuum theory. For example, one can (i) classify
operators under distinct representations of $SU_L(n_f) \times SU_R(n_f)$, and chiral
symmetry prohibits mixing between operators in different representations,
(ii) derive relations between matrix elements of a given operator between
states with different number of pions in the final state (for example
$K \to \pi\pi \leftrightarrow K \to \pi\pi\pi$), (iii) exploit the operator Ward identity, and (iv)
derive relations between correlation functions, i.e. the Ward-Takahashi
identities. It is therefore desirable to preserve chiral symmetry at finite
$a$ and not just in the continuum limit. To determine if this is possible one
needs to understand the relationship between doublers and chiral symmetry
in the lattice theory.

The chiral symmetry defined in Eq.~7.11 is realized by the lattice theory
provided $\gamma_5 M + M \gamma_5 = 0$. It is easy to verify that the na\"ive lattice action
satisfies the hermiticity property
\begin{equation}
\gamma_5 M \gamma_5 = M^\dagger
\end{equation}
Thus, in the massless limit, the Euclidean lattice action has to be
anti-hermitian for chiral symmetry, Eq.~7.11, to hold. This is true of $M^N$
given in Eq.~7.7, and was accomplished by simply taking the symmetric
difference for the derivative. The important question is --- what are the
general conditions under which a lattice theory with this realization of
chiral symmetry exhibits doublers. The answer is given by the Nielsen-
Ninomiya theorem [14]. Consider a generalized action such that
$S^{-1}(p, m = 0) = iF(p)$. Then, if
\begin{itemize}
\item the function $F(p)$ is periodic in momentum space with period $2\pi/a$.
  (This is a consequence of translation invariance.)
\item the lattice momenta are continuous in the range $\{0, 2\pi\}$. This is true in
  the $L \to \infty$ limit at all values of the lattice spacing.
\item $F(p)$ is continuous in momentum space. This is guaranteed if the
  interactions are local.
\item $F(p) \to p_\mu \gamma_\mu$ for small $p_\mu$ and as $a \to 0$ to match the continuum theory.
\item the action has chiral symmetry, i.e. it satisfies the hermiticity property
  $\gamma_5 M \gamma_5 = M^\dagger$ and the massless Dirac action is anti-hermitian.
\end{itemize}
the theory will have doublers. It is easy to understand this from the
schematic one dimensional plot of $F(p)$ shown in Fig.~6 which highlights
the fact that $F(p)$ has to have an even number of zeros under the above
conditions.

\begin{figure}[htbp]
\centering
\includegraphics[width=0.8\textwidth]{images/fig6.png}
\caption{A schematic 1-d representation of the massless inverse propagator $F(p)$.}
\label{fig:6}
\end{figure}

The $\pm$ slopes imply that one state is left moving while the other
is right moving. In 1-dimension the direction is the same as handedness,
consequently left-handed and right-handed fields come in inseparable pairs,
and together form a Dirac fermion.

The key consequence of the Nielsen-Ninomiya theorem is that one cannot
define a lattice theory with chiral fermions. For example if one writes
down $M^N$ for a left-handed Dirac fermion, then its doubler will be
right-handed but have the same transformation under $L \otimes R$. Consequently,
the two ``states'' coexist to form a Dirac fermion. Assuming that the
Nielsen-Ninomiya theorem is air-tight, then how does one proceed to (i)
define a lattice theory without doublers and which recovers chiral symmetry
in the continuum limit, and (ii) define chiral gauge theories on the lattice?

The answer to the first question will be given by explicit reconstruction of
Wilson's and staggered fixes. An alternate very exciting new development,
which I will not do justice to, is the resurrection [53] of the Ginsparg-Wilson
condition for a chirally symmetric formulation [54]
\begin{equation}
\gamma_5 D + D\gamma_5 = a D \gamma_5 D .
\end{equation}
L\"uscher has recently presented a clear analysis of how this condition
allows the construction of a lattice theory without doublers and with chiral
symmetry, avoids the NN theorem, and gives the ABJ anomaly [55]. Un-
fortunately, the bottleneck with this alternative is the lack of a closed form
expression for the lattice Dirac action that is sufficiently local to make
simulations feasible. The second question --- construction of chiral gauge
theories --- is still unresolved. For those interested in this very important
unsolved problem, a good starting point are the reviews [15--17] which
contain some of the recent efforts.
